% Einheit 5 von 5 – Anwendungen (bank/lineare-funktionen/e5.jsonl, zusammenbau v0.1)
%% TODO Zweigzeile Teil 1: was hier gelernt wird (Satz in Schülersprache aus dem Einheitstitel) – Einheit 5
\einheitenkopf[e5]{Einheit 5 von 5 · Anwendungen}
\zweigzeile{ab Kl. 8 · P10 oft}
\setcounter{aufgabe}{27}

%% TODO Titel: Ich-kann-Satz fehlt in der Bank (Platzhalter: Kettenname) – Nr. 28
%% TODO Anweisung: ein Satz über der ersten Teilaufgabe fehlt in der Bank – Nr. 28
\begin{aufgabe}{Anwendung}
\begin{teile}
\teil Markiere: Unterstreiche den Anfangswert, kreise die Änderung je Einheit ein. Ein Handwerker berechnet 45 € für die Anfahrt und 38 € je Arbeitsstunde.
\end{teile}
\begin{teile}
\teil Welche Gleichung passt? Ein Fahrradverleih nimmt 8 € Grundgebühr und 3 € je Stunde. y ist der Preis in €, x die Anzahl der Einheiten. Kreuze an. \\ \kreuz{$y = 8x + 3$} \\ \kreuz{$y = 3x + 8$} \\ \kreuz{$y = 11x$} \\ \kreuz{$y = 3x - 8$}
\end{teile}
\begin{teile}
\teil Welche Gleichung passt? Ein Stromtarif kostet 9 € Grundpreis im Monat und 0{,}35 € je kWh. y ist der Preis in €, x die Anzahl der Einheiten. Kreuze an. \\ \kreuz{$y = 9x + 0{,}35$} \\ \kreuz{$y = 35x + 9$} \\ \kreuz{$y = 0{,}35x + 9$} \\ \kreuz{$y = 0{,}35x - 9$}
\end{teile}
\begin{teile}
\teil Welche Gleichung passt? Ein Taxi kostet 4 € Grundgebühr und 2 € je km. y ist der Preis in €, x die Anzahl der Einheiten. Kreuze an. \\ \kreuz{$y = 4x + 2$} \\ \kreuz{$y = 6x$} \\ \kreuz{$y = 2x - 4$} \\ \kreuz{$y = 2x + 4$}
\end{teile}
\begin{teile}
\teil Welche Gleichung passt? Eine Musikschule verlangt 15 € Anmeldegebühr und 40 € je Monat. y ist der Preis in €, x die Anzahl der Einheiten. Kreuze an. \\ \kreuz{$y = 40x + 15$} \\ \kreuz{$y = 15x + 40$} \\ \kreuz{$y = 55x$} \\ \kreuz{$y = 40x - 15$}
\end{teile}
\begin{teile}
\teil Welche Gleichung passt? Carsharing kostet 1 € Buchungsgebühr und 0{,}29 € je Minute. y ist der Preis in €, x die Anzahl der Einheiten. Kreuze an. \\ \kreuz{$y = x + 0{,}29$} \\ \kreuz{$y = 29x + 1$} \\ \kreuz{$y = 1{,}29x$} \\ \kreuz{$y = 0{,}29x + 1$}
\end{teile}
\swfrage{Was bleibt gleich, was ändert sich?}
%% TODO Darstellung für schwach fehlt in der Bank (lineare-funktionen-e5-k1-s2-v1; Abschnitt „Für schwache Schüler“)
\swz{Stelle eine Gleichung auf: Ein Kino-Abo kostet einmalig 10 € und 6 € je Film. y sind die Kosten in €, x die Anzahl der Filme.}{}
%% TODO Darstellung für schwach fehlt in der Bank (lineare-funktionen-e5-k1-s3-v1; Abschnitt „Für schwache Schüler“)
\swz{Wie viel nach 12 Wochen? Tom hat 85 € gespart. Jede Woche kommen 15 € dazu.}{}
%% TODO Darstellung für schwach fehlt in der Bank (lineare-funktionen-e5-k1-s4-v1; Abschnitt „Für schwache Schüler“)
\swz[3]{Welcher Tarif ist für 60 km günstiger? Tarif A: 20 € Grundgebühr und 0{,}30 € je km. Tarif B: 0{,}70 € je km ohne Grundgebühr.}{}
\swz[3]{$\text{Welches Mietauto? Angebot 1: 32 € Miete pro Tag, 300 km insgesamt frei, jeder weitere km 0{,}30 €. Angebot 2: einmalig 95 € für eine Woche, jeder km 0{,}12 €. Jana mietet 4 Tage und fährt 420 km. Berechne die Gesamtkosten beider Angebote, entscheide, welches günstiger ist, und stelle für Angebot 2 eine Gleichung für die Kosten y (in €) bei x km für eine Woche auf. \hfill (P10 2023 OS)}$}{}
\end{aufgabe}

%% TODO Titel: Ich-kann-Satz fehlt in der Bank (Platzhalter: Kettenname) – Nr. 29
%% TODO Anweisung: ein Satz über der ersten Teilaufgabe fehlt in der Bank – Nr. 29
\begin{aufgabe}{Graph zu Tarif zuordnen}
\swz[3]{Welcher Graph passt zum Tarif? Tarif: 10 € Grundgebühr und 2 € je Stunde. Gib den passenden Graphen an.}{\begin{ksys}[xmin=0,xmax=8,ymin=0,ymax=40,xstep=1,ystep=5,ablesen] \gerade{2}{0}{A} \gerade{2}{10}{B} \gerade{10}{2}{C} \end{ksys}}
\end{aufgabe}

%% TODO Titel: Ich-kann-Satz fehlt in der Bank (Platzhalter: Kettenname) – Nr. 30
%% TODO Anweisung: ein Satz über der ersten Teilaufgabe fehlt in der Bank – Nr. 30
\begin{aufgabe}{Situation zu Gleichung beschreiben}
%% TODO Darstellung für schwach fehlt in der Bank (lineare-funktionen-e5-k3-s1-v1; Abschnitt „Für schwache Schüler“)
\swz[3]{Welche Situation passt? Beschreibe eine Situation zur Gleichung $y = 2{,}5x + 4$ und sage, was x und y bedeuten.}{}
\end{aufgabe}

%% TODO Titel: Ich-kann-Satz fehlt in der Bank (Platzhalter: Kettenname) – Nr. 31
%% TODO Anweisung: ein Satz über der ersten Teilaufgabe fehlt in der Bank – Nr. 31
\begin{aufgabe}{Anwendung – Fehler finden, Begründen, Darstellungswechsel, Anwendung}
\swz[3]{Ich finde den Fehler: Ein Verleih kostet 6 € Grundgebühr und 2{,}50 € je Stunde. Max schreibt $y = 6x + 2{,}5$. Benenne den Fehler und gib die richtige Gleichung an.}{}
\swfrage{Welcher Tarif ist immer günstiger? Tarif A: 10 € Grundgebühr und 0{,}20 € je Minute. Tarif B: 10 € Grundgebühr und 0{,}15 € je Minute. Begründe ohne Rechnung.}
\swz{Gleichung aus der Tabelle: Die Tabelle zeigt die Kosten y (in €) eines Handytarifs für x GB. Gib die Gleichung des Tarifs an.}{\wertetabelle[12,15,18,21]{x}{y}{0,2,4,6}}
\swz[3]{$\text{Wie weit reicht das Geld? Für die Klassenfahrt kostet ein Bus 180 € Grundpreis und 1{,}90 € je km. Die Klasse hat 600 € für den Bus. Wie viele Kilometer (hin und zurück) sind höchstens drin?}$}{}
\end{aufgabe}

\uebersichtskasten{Sachverhalt → Gleichung: n = Anfangswert (Grundgebühr), m = Änderung je Einheit (Preis je km). \\ 3,50 € Grundgebühr, 2 € je km: K(x) = 2x + 3,5}
